\begin{definition} \label{definition:divisibility_of_elements_of_rings}
    Let $R$ be a \CrefAndHyperrefIfExist{definition:ring}{ring} and let $a, b \in R$. We define the following notions of divisibility:
    \begin{enumerate}
        \item We say $a$ \hl{left divides} $b$ (or $a$ is a \hldef{left divisor of $b$}), written \hl{$a \mid_L b$}, if there exists an element $x \in R$ such that $ax = b$.
        \item We say $a$ \hl{right divides} $b$ (or $a$ is a \hldef{right divisor of $b$}), written \hl{$a \mid_R b$}, if there exists an element $y \in R$ such that $ya = b$.
        \item If $R$ is \CrefAndHyperrefIfExist{definition:commutative_ring}{commutative}, these notions coincide, and we simply say \hldef{$a$ divides $b$} or that \hldef{$a$ is a divisor of $b$}, written \hl{$a \mid b$}.
    \end{enumerate}
    In terms of \CrefAndHyperrefIfExist{definition:ideal_generated_by_a_subset_of_a_ring}{principal ideals}, $a \mid_L b$ is equivalent to the containment of \CrefAndHyperrefIfExist{definition:ideal_of_a_ring}{right ideals} $bR \subseteq aR$, and $a \mid_R b$ is equivalent to the containment of \CrefAndHyperrefIfExist{definition:ideal_of_a_ring}{left ideals} $Rb \subseteq Ra$.
\end{definition}